% =====================================================================
% 重写稿（2026-09-28）
% - 引言、相关工作、诊断框架、实验设置：已按核实过的事实重写，可直接使用。
% - 5.1（E1）：使用 e1_ltsf_sampled.csv 的数字；补上置信区间后更新。
% - 5.2（E2）：使用 e2_recomputed.csv（Weather，按"先求和再求比值"重算的 R）；
%   E2 重跑（固定分组、64 个目标、≥150 个窗口）后整体替换。
% - 5.3（E3）、5.4（E5）：原结果存在实现错误，重跑前只保留结构，不写结论。
% - 参考文献使用 references_fixed.bib（key 与原稿一致）。
% =====================================================================
\documentclass[runningheads]{llncs}

\usepackage[T1]{fontenc}
\usepackage{lmodern} % 与 Computer Modern 同一字形的可缩放版本，避免 T1 位图字体
\usepackage{microtype}
\usepackage{graphicx}
\usepackage{booktabs}
\usepackage{amsmath,amssymb}
\usepackage{xcolor}
\usepackage[caption=false]{subfig}
\usepackage[pdfauthor={},hidelinks]{hyperref}

\newif\ifshowtodos
\showtodostrue
\newcommand{\todo}[1]{\ifshowtodos\textcolor{red}{[TODO: #1]}\fi}

\newcommand{\MV}{\textsc{mv}}
\newcommand{\UV}{\textsc{uv}}
\newcommand{\Ben}{B}
\newcommand{\Rel}{R}
\newcommand{\Sen}{S}
\newcommand{\GC}{\mathrm{GC}}
\newcommand{\tfm}{TimesFM-3}

\begin{document}

\title{Does Variate Attention Pay Off? Probing Cross-Channel Use in
Multivariate Time-Series Foundation Models}
\titlerunning{Probing Cross-Channel Use in Multivariate TSFMs}
\author{Anonymous Author(s)}
\authorrunning{Anonymous}
\institute{Anonymous Institution}
\maketitle

\begin{abstract}
% 暂定稿，E2/E3/E5 重跑后更新数字；控制在 200 词以内。
Multivariate time-series foundation models (TSFMs) such as \tfm{} let every
channel attend to the others and report better leaderboard ranks than their
univariate counterparts. Whether these gains come from reading other channels
is unknown. We probe \tfm{} from the inside: we switch its variate attention
between univariate and multivariate modes on identical windows, corrupt the
context channels in controlled ways while keeping the target intact, and
compare the model with Kalman-filter oracles on synthetic systems whose
cross-channel information is known. On eight standard benchmarks the
multivariate mode changes the error by between $-19.5\%$ and $+6.6\%$
depending on the dataset, and on Solar-Energy the loss grows with the
horizon. On Weather, most of the gain survives when the context channels are
replaced by misaligned segments of the same channels, but it disappears, and
turns into a loss, when they are replaced by noise; a leading-indicator
channel, in contrast, is exploited. \todo{E3 and E5 findings.} We release a
diagnostic protocol and derive guidelines for grouping channels when TSFMs
are deployed on high-dimensional sensor data.
\keywords{Time series forecasting \and Foundation models \and
Channel dependence \and Multivariate forecasting \and Model probing}
\end{abstract}

% =====================================================================
\section{Introduction}\label{sec:intro}

Time-series foundation models (TSFMs) forecast unseen series without
task-specific training~\cite{das2024timesfm,ansari2024chronos,woo2024moirai}.
Most of them began as univariate forecasters, and the newest generation adds
attention across channels. \tfm{}~\cite{timesfm3_blog2026} inserts a full
variate-attention sublayer into each of its 20 layers and reports that its
multivariate mode improves the average rank on three public benchmarks.
Chronos-2~\cite{ansari2025chronos2} shares information within groups of
series through group attention.

The evidence that cross-channel attention helps is less clear than these
reports suggest. Moirai~\cite{woo2024moirai} attended across all variates,
yet its successor, Moirai 2.0~\cite{liu2025moirai2}, forecasts
each variate independently. Liu et al.~\cite{liu2026falconx} observe that Chronos-2
performs almost identically on GIFT-Eval with and without group attention.
For supervised models, channel-independent training has long been hard to
beat~\cite{nie2023patchtst,zeng2023dlinear}: Han et
al.~\cite{han2024capacity} attribute this to a trade-off between capacity
and robustness under distribution drift, and Abdelmalak et
al.~\cite{abdelmalak2025lookback} show that standard benchmarks contain weak
inter-channel dependence. Leaderboard gains of a multivariate TSFM are
therefore compatible with two different explanations: the model reads useful
information from other channels, or it benefits from them in a cruder way,
for example by pooling shared seasonal patterns.

Aggregate scores cannot separate these explanations, because they compare
different models or average over heterogeneous tasks. We instead probe a
single model from the inside. Keeping the target channel and the weights
fixed, we switch variate attention between univariate (\UV) and multivariate
(\MV) modes, intervene on the context channels, and measure how accuracy and
the forecast itself respond. We ask four questions:
\textbf{RQ1} Where does the \MV{} mode help, and where does it hurt?
\textbf{RQ2} Which properties of the context channels does the model rely on?
\textbf{RQ3} On systems with known cross-channel information, how much of it
does the model extract?
\textbf{RQ4} How many channels, and which ones, should be supplied when a
dataset has hundreds of them?

Our contributions are as follows.
\begin{itemize}
  \item A within-model diagnostic protocol with four measures (benefit,
        reliance, sensitivity and oracle gap closure) and a set of
        interventions with positive and negative controls
        (Sect.~\ref{sec:framework}).
  \item An analysis of how \tfm{} actually processes more than 32 channels:
        its reference implementation splits them into fixed blocks of 32 in
        column order, so the ``all channels attend to all channels'' setting
        does not occur on high-dimensional benchmarks
        (Sect.~\ref{sec:setup}).
  \item Empirical findings on eight benchmarks \todo{and synthetic systems,
        and a comparison with Chronos-2 and supervised baselines}
        (Sect.~\ref{sec:results}).
  \item Guidelines for selecting channels and deciding between the two modes
        when TSFMs are deployed on sensor data (Sect.~\ref{sec:discussion}).
\end{itemize}

% =====================================================================
\section{Related Work}\label{sec:related}

\subsubsection{Multivariate TSFMs.}
Cross-channel mechanisms in TSFMs differ in scope and cost.
Moirai~\cite{woo2024moirai} flattens all variates into one sequence and
distinguishes them with learned binary attention biases, so its cost grows
quadratically with the number of variates times the sequence length.
Chronos-2~\cite{ansari2025chronos2} applies group attention at each patch
index and learns multivariate behavior from synthetic data in which
dependencies are imposed on univariate series. \tfm{}~\cite{timesfm3_blog2026}
and Toto~2.0~\cite{garg2026toto2} alternate causal attention along time with
full attention across variates, whereas TiRex-2~\cite{carlsen2026tirex2} lets
targets read covariates but not the reverse. Sun et al.~\cite{rmisc2026}
report that replacing synthetic multivariate pretraining data with real
multivariate data improves generalization, which suggests that synthetic data
alone may teach limited cross-channel skill. The closest work to ours is
Falcon-X~\cite{liu2026falconx}, whose authors observe that Chronos-2 is
nearly insensitive to disabling group attention and that its attention maps
are similar for dissimilar inputs. They use this observation to motivate a new
architecture; we turn it into a systematic diagnosis, apply it to \tfm, and
characterize when variate attention helps, is idle, or hurts.

\subsubsection{Channel independence and dependence.}
Channel-independent models, e.g., PatchTST~\cite{nie2023patchtst} and
DLinear~\cite{zeng2023dlinear}, match or outperform many channel-dependent
Transformers~\cite{zhou2021informer,wu2021autoformer,zhou2022fedformer},
while iTransformer~\cite{liu2024itransformer} treats each variate as a token.
Han et al.~\cite{han2024capacity} explain the advantage of channel
independence as robustness to drift. Abdelmalak et
al.~\cite{abdelmalak2025lookback} show that it is partly an artifact of
untuned lookback windows and of weak dependence in standard datasets:
channel-independent models win 21 of 28 comparisons on standard benchmarks,
but channel-dependent models win 18 of 24 on chaotic ODE systems. Song et
al.~\cite{song2026cloc} select a channel-communication budget per dataset on
validation data and find that suppressing communication is preferable on ETT,
whereas local and global communication help on Electricity and Solar-Energy.
Leading indicators are a major source of cross-channel
information~\cite{lift2024}, and channel-dependent models can overfit the
channel order~\cite{cpiri2026}. These studies train a model per dataset; we
ask the corresponding questions for one pretrained model at inference time,
where the use of other channels can be switched off and manipulated without
retraining.

\subsubsection{Evaluating and probing TSFMs.}
Benchmarks such as GIFT-Eval~\cite{gifteval2024},
fev-bench~\cite{fevbench2025} and TIME~\cite{timebench2026} compare models
across many tasks, and fev-bench reports that natively multivariate models
lead on multivariate tasks. Such comparisons confound the effect of
cross-channel attention with every other difference between models. Because
attention weights need not explain predictions~\cite{jain2019attention}, we
intervene on inputs rather than read attention maps, and use surrogate
series~\cite{schreiber1996surrogate} to destroy cross-channel phase relations
while preserving each channel's spectrum and marginal distribution.

% =====================================================================
\section{Diagnostic Framework}\label{sec:framework}

\subsection{Modes and Interventions}
For a test window $w$ and target channel $i$, the model receives the target
history and the histories of a set of context channels $\mathcal{C}_i$ and
returns a forecast of length $H$. In \UV{} mode $\mathcal{C}_i=\emptyset$;
in \MV{} mode $\mathcal{C}_i$ contains the channels processed jointly with
$i$. An intervention $\pi$ modifies only the channels in $\mathcal{C}_i$ and
never the target. We write $\hat y^{\UV}$, $\hat y^{\MV}$ and $\hat y^{\pi}$
for the resulting forecasts and $e(\cdot)$ for their mean squared error on the
standardized scale.

\subsection{Measures}
\emph{Benefit} measures whether the \MV{} mode helps:
\begin{equation}
  \Ben = \frac{\sum_{w,i} e(\hat y^{\UV}_{w,i}) - \sum_{w,i} e(\hat y^{\MV}_{w,i})}
              {\sum_{w,i} e(\hat y^{\UV}_{w,i})}.
\end{equation}
\emph{Reliance} measures how much of the benefit an intervention removes:
\begin{equation}
  \Rel_\pi = \frac{\sum_{w,i} e(\hat y^{\pi}_{w,i}) - \sum_{w,i} e(\hat y^{\MV}_{w,i})}
                  {\sum_{w,i} e(\hat y^{\UV}_{w,i}) - \sum_{w,i} e(\hat y^{\MV}_{w,i})}.
\end{equation}
$\Rel_\pi\approx 1$ means that the benefit depends on what $\pi$ destroys,
$\Rel_\pi\approx 0$ means that it does not, and $\Rel_\pi>1$ means that the
corrupted channels make the forecast worse than using no context at all.
Errors are summed before the ratio is taken, because per-window ratios are
unstable when the two modes perform almost equally. We report $\Rel_\pi$
only when the confidence interval of $\Ben$ excludes zero; otherwise we
report the error change directly. \emph{Sensitivity}
$\Sen_\pi=\operatorname{mean}_{w,i,h}|\hat y^{\pi}-\hat y^{\MV}|/\sigma_i$
measures how much the forecast moves, regardless of accuracy, where
$\sigma_i$ is the training-set standard deviation of channel $i$.
On synthetic linear Gaussian systems, a Kalman filter with the true parameters
gives the lowest achievable error when observing only the target
($e^{\UV}_\star$) or all channels ($e^{\MV}_\star$); \emph{gap closure}
$\GC=(e^{\UV}-e^{\MV})/(e^{\UV}_\star-e^{\MV}_\star)$ is the share of this
available improvement that the model realizes.

\subsection{Interventions and Controls}
Table~\ref{tab:interventions} lists the interventions. The negative controls
must leave the forecast (N1) or its accuracy (N2, N3) unchanged, and the
positive control P1 must improve it; together they validate the protocol
before the remaining interventions are interpreted.

\begin{table}[t]
  \centering
  \caption{Interventions on the context channels $\mathcal{C}_i$. The target
  channel is never modified.}
  \label{tab:interventions}
  \setlength{\tabcolsep}{4pt}
  \begin{tabular}{@{}lll@{}}
    \toprule
    ID & Intervention & What it removes \\
    \midrule
    N1 & permute channels within $\mathcal{C}_i$ & nothing \\
    N2 & add a copy of the target               & nothing new \\
    N3 & add the target lagged by 16 steps      & nothing new \\
    N4 & replace by Gaussian noise, matched moments & all information \\
    I1 & shuffle time steps in each channel & temporal structure \\
    I2 & same channels $\geq$30 days earlier, unaligned & alignment, calendar phase \\
    I3 & as I2, shifted by whole weeks & alignment only \\
    I4 & IAAFT surrogate of each channel & cross-channel phase \\
    I6 & circular shift by 24 steps              & fine alignment \\
    P1 & add $y_i(t+k)+\varepsilon$, $k\in\{16,32\}$ & --- (positive control) \\
    \bottomrule
  \end{tabular}
\end{table}

\subsection{Synthetic Systems}
\todo{Define the generators (low-rank factor, local blocks, structural drift,
lead--lag, independent channels) and the Kalman oracles after the oracle has
passed its unit tests (multivariate oracle error $\leq$ univariate oracle
error; empirical error matches the Riccati steady state).}

% =====================================================================
\section{Experimental Setup}\label{sec:setup}

\subsubsection{Model.}
We use the public \tfm{} checkpoint (330M parameters) and its PyTorch
reference implementation\footnote{Repository commit \texttt{e31dadd}
(2026-09-15).}. Each of its 20 layers applies causal attention along time with
rotary position embeddings, then attention across variates without any
position encoding, then a feed-forward block. The horizon is appended as
masked patches and forecast in a single forward pass. Because the variate
attention carries no position information, it is equivariant to channel
permutations, which N1 verifies.

\subsubsection{How \tfm{} handles more than 32 channels.}
The reference evaluator places at most 32 variates in one forward pass, which
matches the \texttt{max\_variates} value in the released configuration. When
a dataset has more channels, the evaluator splits them into consecutive blocks
of 32 in column order and pads the last block by repeating the first channels.
Traffic (862 channels) is thus processed as 27 independent blocks, and each
channel attends only to the other members of its block. Unless stated
otherwise, our \MV{} results follow this reference protocol.

\subsubsection{Data.}
We use ETTh1, ETTh2, ETTm1, ETTm2 (7 channels), Weather (21),
Solar-Energy (137), Electricity (321) and Traffic (862), with the standard
chronological splits (12/4/4 months for ETT, 7:1:2 otherwise), standardization
by training-set statistics, a context of 512 steps, and
horizons of 96, 192, 336 and 720 steps. For each dataset we evaluate 300 test windows
spaced evenly over the test period. Following the public pretraining
inventory~\cite{gifteval2024}, the pretraining data likely contain the same
862 Traffic sensors, so Traffic is not a clean zero-shot test; the
inventory contains none of the ETT datasets.

\subsubsection{Protocol.}
All runs use FP32 on one RTX 4090 with the evaluator's default symmetric
averaging; we disable clipping of forecasts at zero, which would bias
standardized data. \UV{} and \MV{} forecasts are computed on the same windows,
so all comparisons are paired. \todo{95\% confidence intervals from 2{,}000
bootstrap resamples of windows; Wilcoxon signed-rank tests with Holm
correction.}

% =====================================================================
\section{Results}\label{sec:results}

\subsection{RQ1: Where Does the Multivariate Mode Help?}\label{sec:rq1}

\begin{table}[t]
  \centering
  \caption{\tfm{} in \UV{} and \MV{} mode. MSE is averaged over the four
  horizons; $\Ben$ (\%) is given for the average and for each horizon.
  Positive $\Ben$ means that the \MV{} mode is better. 300 test windows per
  dataset. \todo{add 95\% CIs.}}
  \label{tab:ltsf}
  \setlength{\tabcolsep}{3.6pt}
  \begin{tabular}{@{}lrccrrrrr@{}}
    \toprule
    & & \multicolumn{2}{c}{MSE} & & \multicolumn{4}{c}{$\Ben$ (\%) by horizon} \\
    \cmidrule(lr){3-4}\cmidrule(l){6-9}
    Dataset & $N$ & \UV & \MV & $\Ben$ (\%) & 96 & 192 & 336 & 720 \\
    \midrule
    ETTh1       &   7 & 0.404 & 0.398 &  $+1.6$ &  $+0.1$ &  $+1.1$ &  $+3.6$ &  $+1.3$ \\
    ETTh2       &   7 & 0.329 & 0.307 &  $+6.6$ &  $+6.2$ &  $+6.9$ &  $+7.2$ &  $+6.0$ \\
    ETTm1       &   7 & 0.362 & 0.368 &  $-1.7$ &  $+0.3$ &  $-1.4$ &  $-3.1$ &  $-1.9$ \\
    ETTm2       &   7 & 0.299 & 0.299 &  $0.0$  &  $-0.9$ &  $-0.1$ &  $+0.8$ &  $-0.2$ \\
    Weather     &  21 & 0.301 & 0.287 &  $+4.5$ &  $+8.0$ &  $+3.1$ &  $+1.3$ &  $+5.8$ \\
    Solar-Energy& 137 & 0.256 & 0.306 & $-19.5$ &  $-7.7$ & $-12.7$ & $-20.6$ & $-32.0$ \\
    Electricity & 321 & 0.153 & 0.150 &  $+2.5$ &  $+2.6$ &  $+2.2$ &  $+2.1$ &  $+2.9$ \\
    Traffic$^\dagger$ & 862 & 0.378 & 0.374 & $+1.0$ & $+1.2$ & $+0.6$ & $+0.9$ & $+1.4$ \\
    \bottomrule
    \multicolumn{9}{@{}l}{\footnotesize $^\dagger$Likely seen during pretraining.}
  \end{tabular}
\end{table}

\noindent\textbf{Finding 1. The \MV{} mode helps on some datasets, is neutral
on others, and hurts badly on Solar-Energy.}
Table~\ref{tab:ltsf} shows gains of $+6.6\%$ on ETTh2, $+4.5\%$ on Weather and
$+2.5\%$ on Electricity, no effect on ETTm2, a small loss on ETTm1, and a loss
of $19.5\%$ on Solar-Energy that grows from $7.7\%$ at $H=96$ to $32.0\%$ at
$H=720$. The pattern only partly follows the per-dataset budgets selected for
supervised models: Song et al.~\cite{song2026cloc} suppress channel
communication on all four ETT datasets, whereas \tfm{} gains on ETTh2. A
zero-shot model and a model trained per dataset therefore need not agree on
which datasets reward cross-channel modeling.

\noindent\textbf{The Solar-Energy loss is concentrated in a minority of
windows.} At $H=96$ the median window is $7.6\%$ \emph{better} in \MV{} mode,
while the mean is $7.7\%$ worse; at $H=720$ the median change is $-18.1\%$.
\todo{Relate per-window losses to time of day and to the variance of the
context channels (night-time channels are constant at zero).}

\subsection{RQ2: What Does the Model Rely On?}\label{sec:rq2}
% 数字来自 e2_recomputed.csv（Weather）；E2 重跑后整体替换。

\begin{table}[t]
  \centering
  \caption{Interventions on Weather, where the \MV{} mode reduces MSE by
  $\Ben=12.4\%$ [3.1, 24.1] on the probed windows. $\Delta$MSE is relative to
  the \MV{} mode; brackets are 95\% bootstrap intervals over windows.}
  \label{tab:e2}
  \setlength{\tabcolsep}{5pt}
  \begin{tabular}{@{}llrr@{}}
    \toprule
    ID & Intervention & $\Delta$MSE (\%) & $\Rel_\pi$ \\
    \midrule
    N1 & permute channels         &   $0.0$ & $0.00$ \\
    N2 & copy of target           &  $+1.0$ & $0.07$ [$-0.06$, $0.16$] \\
    N3 & lagged copy of target    &  $+0.6$ & $0.04$ [$-0.02$, $0.12$] \\
    I6 & shift by 24 steps        &  $+2.0$ & $0.15$ [$-0.27$, $0.48$] \\
    I2 & unaligned segment        &  $+2.6$ & $0.18$ [$-0.15$, $0.75$] \\
    I4 & IAAFT surrogate          &  $+6.5$ & $0.46$ [$0.13$, $1.30$] \\
    I3 & week-aligned segment     &  $+6.7$ & $0.48$ [$0.11$, $1.31$] \\
    I1 & shuffled time steps      & $+13.6$ & $0.97$ [$0.53$, $3.10$] \\
    N4 & Gaussian noise           & $+20.1$ & $1.42$ [$0.88$, $3.64$] \\
    \midrule
    P1 & leading indicator, $k=16$ & $-10.0$ & $-0.71$ \\
    P1 & leading indicator, $k=32$ & $-18.6$ & $-1.32$ \\
    \bottomrule
  \end{tabular}
\end{table}

\noindent\textbf{Finding 2. On Weather, the gain survives misaligned context
but not unstructured context.}
Table~\ref{tab:e2} orders the interventions by how much of the benefit they
remove. Replacing the context channels with the same channels from another,
unaligned period (I2) or shifting them by 24 steps (I6) removes less than a
fifth of the benefit, whereas shuffling their time steps (I1) removes all of
it and Gaussian noise (N4) makes the forecast worse than using no context at
all. The model therefore benefits mainly from context that looks like
realistic weather signals, and much less from their exact alignment with the
target, which is closer to pooling shared patterns than to reading
synchronized co-movement. \todo{I2 retains more of the benefit than the
week-aligned I3, contrary to expectation; the intervals overlap, so confirm
with the larger rerun before interpreting.}

\noindent\textbf{Finding 3. The model can exploit a genuine leading
indicator, and cannot ignore noise.}
The negative controls behave as required: permutation leaves the forecast
unchanged, and redundant copies of the target barely change the error. A
channel that reveals the target 16 or 32 steps ahead lowers the error by
$10.0\%$ and $18.6\%$, so the variate attention can route predictive
information when it exists. Noise channels, however, raise the error by
$20.1\%$: the model does not learn to disregard uninformative context.
\todo{Report Electricity, Solar and Traffic after the rerun with fixed
groups; report $\Delta$MSE instead of $\Rel$ where the interval of $\Ben$
contains zero.}

\subsection{RQ3: Synthetic Systems with Known Cross-Channel Information}\label{sec:rq3}
\todo{Rerun after the Kalman oracle passes its unit tests. The current results
are not usable: the model beats the multivariate oracle in 30 of 39
configurations, which is impossible for a correct oracle.}

\subsection{RQ4: How Many Channels, and Which?}\label{sec:rq4}
\todo{Rerun with $K>32$ passed directly to the forecaster rather than through
the evaluator's 32-channel blocks; the current $K=64$ and $K=128$ results are
identical to $K=32$ because the model never saw more than 32 channels.
Current data (for $K\leq 32$, to be confirmed): on Electricity, random
selection is slightly better than top-$K$ selection at every $K\geq 4$; on
Solar-Energy, every added channel increases the error for all three
strategies.}

% =====================================================================
\section{Discussion and Guidelines}\label{sec:discussion}
\todo{Write after the reruns. Candidate guidelines that the current data
already support: (1) validate the \MV{} mode per dataset before deployment,
since it can lose up to 32\% (Solar-Energy); (2) never pad channel groups with
noise-like or unrelated series, because the model does not ignore them.
Mode-selection thresholds and channel-budget rules require E5 and E7.}

\subsubsection{Limitations.}
We probe one model in depth; Chronos-2 and supervised baselines are compared
on a reduced protocol \todo{E4}. Traffic is likely part of the pretraining
data. Interventions change the input distribution, so part of their effect may
reflect distribution shift rather than the loss of specific information; the
negative controls bound this effect.

\section{Conclusion}\label{sec:conclusion}
\todo{Write last.}

\bibliographystyle{splncs04}
\bibliography{references_fixed}

\end{document}
